% Stemma: Grundlagendokument (ConTeXt, MkIV)
% Kompilieren im Ordner docs: context main.tex

\setupbodyfont[modern,11pt]
\mainlanguage[de]
\setuppapersize[A4][A4]
\setuplayout[
  width=middle, height=middle,
  header=1.2cm, footer=1.2cm,
  topspace=2.2cm, bottomspace=2.2cm,
  backspace=2.5cm, margin=2cm]
\setupheadertexts[\getmarking[chapter]][\pagenumber]
\setupfootertexts[][\pagenumber]
\setuphead[chapter][before={\blank[2*big]}, after={\blank[big]}]

\starttext

\startalign[center]
\blank[4*big]
{\tfd Stemma}\par
\blank[big]
{\tfa Ein API-first Versionskontrollsystem für Menschen und Agenten}\par
\blank[2*big]
Grundlagendokument, Entwurf\par
\blank[4*big]
Daniel Senften\par
Talent Factory GmbH\par
\blank[big]
Version 0.2 · 18. Oktober 2026\par
\stopalign
\page

\title{Inhaltsverzeichnis}
\placecontent
\page

\input chapters/01-einleitung
\input chapters/02-grundlagen
\input chapters/03-architektur
\input chapters/04-phase1-lokaler-kern
\input chapters/05-phase2-diff
\input chapters/06-phase3-branches-merge
\input chapters/07-phase4-zusammenarbeit
\input chapters/08-phase5-netzwerk
\input chapters/09-phase6-agenten-schnittstellen
\input chapters/10-lehre
\input chapters/11-ausblick

\title{Abbildungsverzeichnis}
Abbildungen werden ergänzt, sobald sie vorliegen.
\placelistoffigures
\page

\title{Literatur}
Die Quellen folgen dem IEEE-Format. Einträge werden ergänzt, sobald die Quellen feststehen. Das Format für einen Eintrag lautet:

\starttyping
[1] A. Autor, „Titel der Quelle,“ Verlag, Ort, Jahr.
\stoptyping

\stoptext
